% GENERATED FILE. Edit canonical lessons, then run npm run generate:books.
% canonical-source-hash: fnv1a64:c6dd00517eccacfd
% canonical-lessons: BN-C483-jamai, BN-C483-bouma, BN-C483-nonod, BN-C483-deor, BN-C483-bhashur

\chapter{In-Laws}
\label{ch:bn-in-laws}
\hlchaptermodality{483}

\begin{chapteropening}
\textbf{By the end of this chapter:} \emph{I can name a son-in-law, a daughter-in-law and a husband's family.}

Complete the last lesson of chapter 483: I can name a son-in-law, a daughter-in-law and a husband's family.
\end{chapteropening}

\section[\texorpdfstring{\bn{জামাই} (jāmāi)}{জামাই (jāmāi)}]{\bn{জামাই} (jāmāi) — a son-in-law}

\label{lesson:BN-C483-jamai}



\begin{quote}
Before the new one: say the Bengali for a brother's daughter, then the Bengali for a child's father-in-law.
\end{quote}

\subsection*{You'll want to know: \bn{জামাই}}
\textbf{\bn{জামাই}} — \emph{jāmāi} — \textquotedblleft{}a son-in-law\textquotedblright{}.

\begin{quote}
\bn{জামাই} \bn{এলে} \bn{শাশুড়ি} \bn{মাছ} \bn{রান্না} \bn{করেন।}
\end{quote}

\begin{quote}
\emph{jāmāi ele shāshuṛi māchh rānnā kôren.}
\end{quote}

\begin{quote}
When the son-in-law visits, his mother-in-law cooks fish.
\end{quote}

\begin{cousinweb}
\textbf{\bn{জামাই}} (\emph{jāmāi}) comes from Sanskrit \emph{jāmātṛ}, \textquotedblleft{}a daughter's husband\textquotedblright{} (CDIAL 5198), which CDIAL lists in this Bengali form.
\end{cousinweb}

\begin{grammarlens}[title={What you've built}]
One more word for describing what you lived and felt.
\end{grammarlens}

\subsection*{Your turn}
\begin{itemize}
\raggedright
  \item \emph{Say it:} \emph{jāmāi}
  \item \emph{Say it:} \emph{jāmāi}, once more
  \item \emph{Say it:} \emph{jāmāi ele shāshuṛi māchh rānnā kôren.} — When the son-in-law visits, his mother-in-law cooks fish.
  \item \emph{Recall:} say the Bengali for a brother's daughter, then the Bengali for a child's father-in-law, then say \emph{jāmāi} again
\end{itemize}

\begin{tcolorbox}[breakable,colback=teal!4,colframe=teal!35!black,title={Before you move on}]
What does \bn{জামাই} mean? (\textquotedblleft{}A son-in-law\textquotedblright{}.) Say it once more.
\end{tcolorbox}
\section[\texorpdfstring{\bn{বউমা} (boumā)}{বউমা (boumā)}]{\bn{বউমা} (boumā) — a daughter-in-law}

\label{lesson:BN-C483-bouma}



\begin{quote}
Before the new one: say the Bengali for a child's father-in-law, then the Bengali for a son-in-law.
\end{quote}

\subsection*{You'll want to know: \bn{বউমা}}
\textbf{\bn{বউমা}} — \emph{boumā} — \textquotedblleft{}a daughter-in-law\textquotedblright{}.

\begin{quote}
\bn{শাশুড়ি} \bn{বউমাকে} \bn{একটা} \bn{শাড়ি} \bn{দিলেন।}
\end{quote}

\begin{quote}
\emph{shāshuṛi boumāke ekṭā shāṛi dilen.}
\end{quote}

\begin{quote}
The mother-in-law gave her daughter-in-law a sari.
\end{quote}

\begin{cousinweb}
\textbf{\bn{বউমা}} (\emph{boumā}) puts \emph{bou}, \textquotedblleft{}wife, bride\textquotedblright{}, from Sanskrit \emph{vadhū} (CDIAL 11250), with \bn{মা}, mother: the bride who becomes a daughter of the house.
\end{cousinweb}

\begin{grammarlens}[title={What you've built}]
One more word for describing what you lived and felt.
\end{grammarlens}

\subsection*{Your turn}
\begin{itemize}
\raggedright
  \item \emph{Say it:} \emph{boumā}
  \item \emph{Say it:} \emph{boumā}, once more
  \item \emph{Say it:} \emph{shāshuṛi boumāke ekṭā shāṛi dilen.} — The mother-in-law gave her daughter-in-law a sari.
  \item \emph{Recall:} say the Bengali for a child's father-in-law, then the Bengali for a son-in-law, then say \emph{boumā} again
\end{itemize}

\begin{tcolorbox}[breakable,colback=teal!4,colframe=teal!35!black,title={Before you move on}]
What does \bn{বউমা} mean? (\textquotedblleft{}A daughter-in-law\textquotedblright{}.) Say it once more.
\end{tcolorbox}
\section[\texorpdfstring{\bn{ননদ} (nônôd)}{ননদ (nônôd)}]{\bn{ননদ} (nônôd) — a husband's sister}

\label{lesson:BN-C483-nonod}



\begin{quote}
Before the new one: say the Bengali for a son-in-law, then the Bengali for a daughter-in-law.
\end{quote}

\subsection*{You'll want to know: \bn{ননদ}}
\textbf{\bn{ননদ}} — \emph{nônôd} — \textquotedblleft{}a husband's sister\textquotedblright{}.

\begin{quote}
\bn{মীরা} \bn{ননদের} \bn{সঙ্গে} \bn{বাজারে} \bn{গেল।}
\end{quote}

\begin{quote}
\emph{mīrā nônôder sôngge bājāre gelô.}
\end{quote}

\begin{quote}
Mira went to the market with her husband's sister.
\end{quote}

\begin{cousinweb}
\textbf{\bn{ননদ}} (\emph{nônôd}) comes from Sanskrit \emph{nanāndṛ}, \textquotedblleft{}a husband's sister\textquotedblright{} (CDIAL 6946), which CDIAL lists in this Bengali form.
\end{cousinweb}

\begin{grammarlens}[title={What you've built}]
One more word for describing what you lived and felt.
\end{grammarlens}

\subsection*{Your turn}
\begin{itemize}
\raggedright
  \item \emph{Say it:} \emph{nônôd}
  \item \emph{Say it:} \emph{nônôd}, once more
  \item \emph{Say it:} \emph{mīrā nônôder sôngge bājāre gelô.} — Mira went to the market with her husband's sister.
  \item \emph{Recall:} say the Bengali for a son-in-law, then the Bengali for a daughter-in-law, then say \emph{nônôd} again
\end{itemize}

\begin{tcolorbox}[breakable,colback=teal!4,colframe=teal!35!black,title={Before you move on}]
What does \bn{ননদ} mean? (\textquotedblleft{}A husband's sister\textquotedblright{}.) Say it once more.
\end{tcolorbox}
\section[\texorpdfstring{\bn{দেওর} (deor)}{দেওর (deor)}]{\bn{দেওর} (deor) — a husband's younger brother}

\label{lesson:BN-C483-deor}



\begin{quote}
Before the new one: say the Bengali for a daughter-in-law, then the Bengali for a husband's sister.
\end{quote}

\subsection*{You'll want to know: \bn{দেওর}}
\textbf{\bn{দেওর}} — \emph{deor} — \textquotedblleft{}a husband's younger brother\textquotedblright{}.

\begin{quote}
\bn{দেওর} \bn{বউদির} \bn{সঙ্গে} \bn{ঠাট্টা} \bn{করে।}
\end{quote}

\begin{quote}
\emph{deor boudir sôngge ṭhāṭṭā kôre.}
\end{quote}

\begin{quote}
A husband's younger brother jokes with his elder brother's wife.
\end{quote}

\begin{cousinweb}
\textbf{\bn{দেওর}} (\emph{deor}) comes from Sanskrit \emph{devṛ}, \textquotedblleft{}a husband's brother\textquotedblright{} (CDIAL 6546), which CDIAL lists in this Bengali form.
\end{cousinweb}

\begin{grammarlens}[title={What you've built}]
One more word for describing what you lived and felt.
\end{grammarlens}

\subsection*{Your turn}
\begin{itemize}
\raggedright
  \item \emph{Say it:} \emph{deor}
  \item \emph{Say it:} \emph{deor}, once more
  \item \emph{Say it:} \emph{deor boudir sôngge ṭhāṭṭā kôre.} — A husband's younger brother jokes with his elder brother's wife.
  \item \emph{Recall:} say the Bengali for a daughter-in-law, then the Bengali for a husband's sister, then say \emph{deor} again
\end{itemize}

\begin{tcolorbox}[breakable,colback=teal!4,colframe=teal!35!black,title={Before you move on}]
What does \bn{দেওর} mean? (\textquotedblleft{}A husband's younger brother\textquotedblright{}.) Say it once more.
\end{tcolorbox}
\section[\texorpdfstring{\bn{ভাশুর} (bhāshur)}{ভাশুর (bhāshur)}]{\bn{ভাশুর} (bhāshur) — a husband's elder brother}

\label{lesson:BN-C483-bhashur}



\begin{quote}
Before the new one: say the Bengali for a husband's sister, then the Bengali for a husband's younger brother.
\end{quote}

\subsection*{You'll want to know: \bn{ভাশুর}}
\textbf{\bn{ভাশুর}} — \emph{bhāshur} — \textquotedblleft{}a husband's elder brother\textquotedblright{}.

\begin{quote}
\bn{বিয়ের} \bn{পরে} \bn{মীরা} \bn{ভাশুরকে} \bn{প্রথম} \bn{দেখল।}
\end{quote}

\begin{quote}
\emph{biyer pôre mīrā bhāshurke prothom dekhlô.}
\end{quote}

\begin{quote}
After the wedding Mira met her husband's elder brother for the first time.
\end{quote}

\begin{cousinweb}
\textbf{\bn{ভাশুর}} (\emph{bhāshur}) is listed by CDIAL for Bengali under Sanskrit \emph{bhrātṛśvaśura}, \textquotedblleft{}brother-father-in-law\textquotedblright{} (CDIAL 9669): an elder brother of the husband, treated with a father-in-law's respect.
\end{cousinweb}

\begin{grammarlens}[title={What you've built}]
One more word for describing what you lived and felt.
\end{grammarlens}

\subsection*{Your turn}
\begin{itemize}
\raggedright
  \item \emph{Say it:} \emph{bhāshur}
  \item \emph{Say it:} \emph{bhāshur}, once more
  \item \emph{Say it:} \emph{biyer pôre mīrā bhāshurke prothom dekhlô.} — After the wedding Mira met her husband's elder brother for the first time.
  \item \emph{Recall:} say the Bengali for a husband's sister, then the Bengali for a husband's younger brother, then say \emph{bhāshur} again
\end{itemize}

\begin{tcolorbox}[breakable,colback=teal!4,colframe=teal!35!black,title={Before you move on}]
What does \bn{ভাশুর} mean? (\textquotedblleft{}A husband's elder brother\textquotedblright{}.) Say it once more.
\end{tcolorbox}
